\documentclass[11pt]{article}
\usepackage[T1]{fontenc}
\usepackage{lmodern,amsmath,amssymb,amsthm,booktabs,microtype}
\usepackage[margin=0.9in]{geometry}
\usepackage[hidelinks]{hyperref}
\newtheorem{lemma}{Lemma}
\newtheorem{theorem}{Theorem}
\newtheorem{corollary}{Corollary}
\newcommand{\bits}{\{0,1\}}
\newcommand{\B}{\mathsf B}
\newcommand{\M}{\mathsf M}
\newcommand{\width}{\operatorname{width}}
\newcommand{\Inc}{\operatorname{inc}}
\newcommand{\Dec}{\operatorname{dec}}
\newcommand{\Con}{\operatorname{Con}}
\newcommand{\concat}{\mathbin{+\!+}}
\title{A 250-state shared-carry machine:\\local identities, runtime transport, and bootstrap}
\author{ProveIt research contribution}
\date{17 October 2026}
\begin{document}
\maketitle

\begin{abstract}
A rightward carry circuit can implement both increment and decrement in the
existing nine-register binary interpreter.  For decrement it shifts the
register prefix, stopping at the selected mark; a zero result rolls that
prefix back.  Increment deposits a harmless spare register of value one.
Retiming a temporary bit makes an otherwise unused read-zero transition a
blank-tape initializer.  The resulting table has 250 working states and 500
binary transitions.  The new register identities are proved below for
arbitrary register values, arbitrary left contexts, and arbitrary nonempty
spare suffixes.  An adapted finite symbolic verifier checks all 8,391 original
instruction guard cases.  A literal 585,431-step bootstrap reaches a safe root.
This is a conventional mathematical proof with executable finite
certificates; it does not claim a newly compiled Lean or Rocq endpoint.
\end{abstract}

\section{Exact machine and conventions}

The source is the frozen original-label permutation252 table in ProveIt,
at commit
\begin{quote}\small\ttfamily
a3f06a2059dd786ad385079106454efb84b87d3c.
\end{quote}
Its SHA-256 is
\begin{quote}\small\ttfamily
939f99d602c8f7057517c4369279de99e30827050fc743581df83f26aaf1b713.
\end{quote}
The final original-label file \texttt{machine/zfc\_prefix250\_old.tm} has SHA-256
\begin{quote}\small\ttfamily
fd455304bad82a6da7ab83149352d0b1db27557582c4cc6312de003c96ff4e43.
\end{quote}
The start state is old 97; old 254 remains the separate halt label.  There
are 250 working labels: the old labels $0,\ldots,253$ except
$5,8,43,59,85$, together with the new label 255.  A canonical renumbering is
also supplied.  The tape is bi-infinite, initially all zero, and the head
starts at coordinate zero.  A table entry $(w,d,q')$ writes $w$, moves one
cell in direction $d$, and enters $q'$.

Delete states 43, 59 and 85 from permutation252 and make precisely these
changes, leaving every other surviving row fixed:
\begin{center}
\begin{tabular}{rrlll}
\toprule
State & Read & Write & Move & Next state\\
\midrule
30 & 1 & 0 & R & 95\\
75 & 0 & 1 & R & 255\\
75 & 1 & 1 & L & 63\\
97 & 0 & 1 & R & 83\\
97 & 1 & 1 & L & 131\\
109 & 0 & 0 & R & 75\\
131 & 0 & 0 & L & 48\\
255 & 0 & 1 & L & 97\\
255 & 1 & 1 & L & 131\\
\bottomrule
\end{tabular}
\end{center}
The state-30 change replaces the old decrement seeker by the carry entry.
No surviving row targets a deleted state.  In
particular, changing both state-97 rows removes its old reference to 85.
The unique transition to halt retains its original write and movement.

\section{Tape notation and exact local contracts}

For a finite list of natural numbers $v=(v_0,\ldots,v_{m-1})$, put
\[
 \B(v)=\prod_{i<m} 0\,1^{v_i+1},\qquad
 \width(v)=\sum_{i<m}(v_i+2).
\]
Thus a physical register block always contains at least one one-cell.  If
$v=a\concat[x]\concat b$, where $|a|=r$, its selected, marked form is
\[
 \M_r(v)=\B(a)\,00\,1^x\B(b).
\]
The first one-cell of register $r$ has been erased.  All cells to the right
of the displayed finite file are zero.  The left tape context $\ell$ is
arbitrary; none of the local contracts imposes PC, marker, or gap conditions
on it.

In each lemma put the first file cell at coordinate zero.  An output ``at
the old left boundary'' means head coordinate $-1$, immediately before the
file, before that cell has been read.  ``At the new left boundary'' means
head coordinate $0$, after the file's boundary has moved one cell right.
The left context is therefore completely preserved in the first case,
and the newly written one at coordinate zero extends it in the second case.
The symbols $\Inc_r(v)$ and $\Dec_r(v)$ change only entry $r$ by $+1$ and
$-1$, respectively; the latter is used only when $v_r>0$.

All nine logical registers occupy physical indices $0,\ldots,8$.  The
global invariant will require $m\geq10$.  Consequently, whenever $r<9$,
the suffix $b$ in the marked decomposition is nonempty.  This is the exact
condition needed to distinguish a zero register from the end of the file.

\subsection{A three-step retiming identity}

For clarity, first consider an auxiliary carry circuit with the same 250
working labels.  Its rows differ from the final table as follows:
\[
\begin{array}{c|cc}
 & \text{read }0 & \text{read }1\\ \hline
75 & (0,R,255)&(1,L,63)\\
255&(1,L,112)&(1,L,97)\\
97&(0,L,112)&\text{unused}
\end{array}
\]
Every other row is already final.  The auxiliary circuit is used only to
prove local identities, not as the delivered blank-start machine.

\begin{lemma}[Retiming]
Start in state 75 scanning a zero at coordinate $c$.  For either value of
the next cell $d=c+1$, the auxiliary and final circuits reach identical
configurations after exactly three steps.  The assertion holds for every
surrounding tape context.
\end{lemma}
\begin{proof}
Only coordinates $c$ and $d$ are written.  If $d$ contains zero, the
auxiliary state sequence is
$75\to255\to112\to131$ and the final sequence is
$75\to255\to97\to131$.
Both end at coordinate $c-1$ in state 131, with both $c$ and $d$ equal to
one.  If $d$ contains one, the auxiliary sequence is
$75\to255\to97\to112$ and the final sequence is
$75\to255\to131\to112$.
Both end at $c-1$ in state 112, with $c=0$ and $d=1$.  These are literal
three-row calculations.  In the final circuit, the temporary one at $c$
ensures that state 97 is entered on read one.  Its read-zero row is free
to initialize the blank tape.
\end{proof}

The retiming preserves both local times and outputs.  We may therefore
derive the remaining contracts with the simpler auxiliary circuit and
transfer them to the exact final table.

\subsection{Positive decrement by moving the prefix}

The unchanged carry rows are
\[
\begin{array}{c|cc}
 &0&1\\ \hline
95&(1,R,109)&(1,R,95)\\
109&(0,R,75)&(0,R,95).
\end{array}
\]
The state-109 read-zero row is the new termination test.  The old return
scanner has
\[
63: (0,L,48),(1,L,63),\qquad
48: (0,L,7),(1,L,63).
\]

\begin{lemma}[Positive selected register]\label{lem:positive}
Let $v=a\concat[x]\concat b$ with $x>0$, and set $p=\width(a)$.
Starting in state 95 at coordinate zero on $\M_r(v)$, the final machine
reaches state 48 at coordinate zero after exactly $2p+4$ transitions.
The tape is
\[
 \ell\;1\;\B(\Dec_r(v))\;0^\infty.
\]
The displayed one is at the original file origin, so the register file
and its header boundary have shifted one cell to the right.  No spare
value or number of spare blocks changes.
\end{lemma}
\begin{proof}
At a normal block $01^{u+1}$, states 95 and 109 replace its initial $01$
by $10$, then state 95 scans the remaining $u$ ones.  Applying this
calculation successively to the blocks of $a$, followed by the first
zero of the mark, takes $p+2$ rightward steps.  The machine is then in
state 75 at coordinate $p+2$, and the tape prefix is
\[
 \ell\;1\,\B(a)\,0\,1^x\B(b).
\]
Since $x>0$, state 75 reads one and enters the return scanner.  Its
isolated-zero block scan takes $p+2$ further leftward steps, stopping
in state 48 at coordinate zero.  The selected block is $01^x$, so its
new value is $x-1$.  Every preceding block merely moved right; every
following block remained in place.  The left context was not read.
\end{proof}

\subsection{Zero detection and rollback}

The reverse carry rows are
\[
\begin{array}{c|cc}
 &0&1\\ \hline
112&(1,L,131)&(1,L,112)\\
131&(0,L,48)&(0,L,112).
\end{array}
\]
The read-zero exit of 131 is deliberately state 48, not its old state 63.
It supplies the zero-decrement control continuation after the local box
has been restored.

\begin{lemma}[Zero selected register]\label{lem:zero}
Let $v=a\concat[0]\concat b$, with $b\ne[]$, and $p=\width(a)$.
Starting in state 95 at coordinate zero on $\M_r(v)$, the final machine
reaches state 112 at coordinate $-1$ after exactly $2p+7$ transitions.
The entire tape is again $\ell\B(v)0^\infty$.
\end{lemma}
\begin{proof}
The first $p+2$ steps are the same carry calculation as in
Lemma~\ref{lem:positive}.  The local suffix is now
\[
 \underbrace{1}_{p}\;
 \underbrace{0}_{p+1}\;
 \underbrace{0}_{p+2}\;
 \underbrace{1}_{p+3}\cdots .
\]
The last one exists because $b$ begins with a valid block.  In the
auxiliary circuit, state 75 at $p+2$ takes the zero branch, 255 sees the
one at $p+3$, and 97 returns left through $p+2$ to state 112 at $p+1$.
The next two transitions restore the selected block: 112 writes one at
$p+1$, and 131 changes the one at $p$ back to zero.  State 112 now scans
the shifted prefix from coordinate $p-1$.

For each preceding block, the pair $112/131$ changes its shifted $10$
back to $01$ while moving left.  Induction on the blocks of $a$ takes
exactly $p$ more steps and restores $\B(a)$ completely.  The result is
state 112 at $-1$, before reading the left context.  The total is
$(p+2)+5+p=2p+7$.  The retiming lemma transfers this exact result to the
final table.
\end{proof}

The output state 112 is part of the contract.  It must not be silently
replaced by the old zero-decrement return state 3.  At an actual
instruction boundary, the subsequent literal $112/131$ scan changes
the two phase bits from $00$ to $01$ before reaching state 48 at the
last PC bit.  The finite instruction verifier executes this change
explicitly and verifies the required break continuation.

\subsection{Increment and the disposable spare}

The unchanged seeker is
$73:(0,R,83),(1,R,73)$ and
$83:(1,R,95),(1,R,73)$.
It finds the selected pair of zeros and restores the erased mark.

\begin{lemma}[Increment]\label{lem:increment}
Let $r<|v|$ and $w=\width(v)$.  Starting in state 73 at coordinate zero
on $\M_r(v)$, the final machine reaches state 48 at coordinate $-1$
after exactly $2w+7$ transitions.  Its tape is
\[
 \ell\;\B\bigl(\Inc_r(v)\concat[1]\bigr)\;0^\infty.
\]
Thus increment appends exactly one spare block, of value one.
\end{lemma}
\begin{proof}
The seeker passes the preceding ordinary blocks, restores the first
one of the selected block, and enters 95.  The carry shifts every later
separator right.  From the file origin to state 75 at coordinate $w+2$
there are exactly $w+2$ rightward transitions.  The last newly written
one is at $w$, and coordinates $w+1,w+2,w+3$ are zero.

In the auxiliary circuit, the zero/zero test writes ones at $w+3$ and
$w+2$, then enters state 48 at $w$.  These four transitions leave the
new spare block $011$ at $w+1,w+2,w+3$.  The return scanner takes $w+1$
further steps to the old left boundary.  The total is
$(w+2)+4+(w+1)=2w+7$.  The carry increased the selected register by one
and preserved all other old values.  Retiming gives the same final
result and duration.
\end{proof}

\subsection{Marking and the strengthened invariant}

The nine selector entry states remain
\[
(60,44,31,56,38,40,28,64,69).
\]
Their previously established marking contract is unchanged: entry $r$
on $\B(v)$ produces $\M_r(v)$ and returns to state 2 at the old left
boundary when $r=0$, and state 6 otherwise.  Every table row used by this
particular excursion is literally unchanged.  The original proof by
induction over skipped unary blocks therefore applies.

\begin{corollary}[Nonempty spare invariant]
Suppose a runtime frame contains nine primary registers and at least
one spare block.  Marking preserves the file.  Increment preserves
all old spare values and appends one value-one spare.  Positive and
zero decrement preserve the spare suffix exactly.  Consequently every
subsequent runtime frame again contains a nonempty spare suffix, and
Lemma~\ref{lem:zero} is available at every physical register $r<9$.
\end{corollary}

The spare values need not be zero or one for the local proofs.  Only
their finite unary-block shape and nonemptiness matter.  The actual
bootstrap happens to provide 210 value-one spares.

\section{All instruction contracts}

A normal dispatch frame has old state 96 and tape
\[
0^\infty\,1\,P\,0^{D-|P|}\,1\,\B(v)\,0^\infty,
\]
with the head at the first PC cell, $D\ge40$, and a nonempty spare suffix.
The unchanged logical-to-physical assignment, in the declaration order
$(N,T,W,A,D,H,X,Y,L)$, is
\[
(6,2,0,4,3,7,8,5,1).
\]
The overloaded letter $D$ in the logical-register list is a register
name; the $D$ in the display is the marker-to-header span.

The supplied \texttt{prefix\_affine.py} adapts the attributed source-27
checker.  The only new symbolic tape operation is to prepend one
one-cell to the left part when applying the positive-decrement
contract.  This is exactly the boundary motion of
Lemma~\ref{lem:positive}.  The abstract contracts are therefore:
\begin{center}
\begin{tabular}{llll}
\toprule
Entry & Condition & Exit state & File / boundary effect\\
\midrule
Selector $r$ & Plain file & 2 or 6 & Mark $r$; fixed boundary\\
73 & Marked $r$ & 48 & Increment $r$; append $[1]$\\
95 & Marked $r$, $v_r=0$ & 112 & Restore file; fixed boundary\\
95 & Marked $r$, $v_r>0$ & 48 & Decrement $r$; boundary $+1$\\
\bottomrule
\end{tabular}
\end{center}
Both decrement guards name the same physical register that was marked.

\begin{lemma}[Finite symbolic verification]
With these four contracts, all 8,391 instruction guard cases of the
unchanged 6,460-leaf decoder satisfy their original register-program
transition, for every additional gap parameter $n\ge0$, every primary
register vector satisfying the named guard, and every nonempty finite
spare suffix.  The final span is
\[
 D'=D+\mathbf1_{\text{positive decrement}}
       +\mathbf1_{\text{PC overflow}}.
\]
The expected register operation and PC successor are unchanged.
\end{lemma}
\begin{proof}
The finite certificate is reproduced by
\texttt{check\_prefix\_affine.py}.  Its inputs are the exact final
transition table and the complete extracted instruction-site list,
whose provenance and hashes are supplied.  Every case starts with
span $40+n$.  A run consists of literal steps, normalization of adjacent
constant or affine runs, whole-run self-loop steps, and the four local
contracts proved above.  A one-cell peel requires a positive constant
part, so no branch on an unknown affine length is assumed.  The extra
one-cell boundary extension has the semantics just described.

Induction over each accepted symbolic trace proves its concrete
meaning for every $n\ge0$ and every allowed file.  The verifier checks
the entire target PC, span and box effect, not merely a return to state
96.  It accepts 8,391 cases; the smallest number of literal transitions
in a case is 13, and at most 497 symbolic rules are used.  The result
records 1,868,409 literal rules over all cases.  Thus all instruction
simulations make positive progress.  These are finite exhaustive
symbolic certificates under proved local identities, rather than an
inference from a bounded machine run.
\end{proof}

The Python checker itself has not been newly verified in a proof
assistant.  Its small transition and symbolic-rule implementation is
part of the reviewable computational proof.  Independent literal tests
are useful regression evidence but are not substituted for the
universal lemmas or symbolic trace argument.

\section{Blank start and global transport}

The final state-97 read-zero row writes one and enters state 83 one
cell to the right.  During the proved runtime excursions, state 97 is
entered only on read one.  This separates the blank initialization from
the runtime without an additional state.

\begin{lemma}[Exact bootstrap]
From the all-zero tape in old state 97 at coordinate zero, the final
machine reaches old state 96 after exactly 585,431 transitions, without
an earlier halt.  The head is at $-3$, the marker at $-4$, the header at
225, and the span is 228.  The PC is empty.  The primary physical vector is
\[
 [0,2,0,4,0,0,1,0,0],
\]
followed by exactly 210 spare blocks, each of value one.  The tape has
exactly 438 one-cells, and every other cell is zero.
\end{lemma}
\begin{proof}
This is a finite literal calculation, reproduced by
\texttt{code/check\_boot250.py}.  Its output records every nonzero
coordinate, and compares the complete final tape with the claimed
frame.  The only visit to state 97 on read zero is the first transition.
The output table hash is the exact frozen hash above.  The corresponding
logical vector is $[1,0,0,0,4,0,0,0,2]$.
\end{proof}

The original frontend's safety and completeness result permits precisely
this kind of residual data: the top formula, formula stack and cached
formula are zero, and the remaining fields are admissible data for its
root loop.  The next-candidate counter is one.  The accepted initial
condition and unchanged logical calibration are part of the accompanying
endpoint audit; no logical axiom, inference rule, or decoder leaf has
been deleted in this construction.

\begin{theorem}[Runtime transport]
Let $\mathcal P$ be the unchanged source register program and let
$R_*=[1,0,0,0,4,0,0,0,2]$ in its declaration order.  The final 250-state
machine halts from the blank tape if and only if $\mathcal P$ halts from
its root with register vector $R_*$.
\end{theorem}
\begin{proof}
The bootstrap establishes the dispatch invariant, including $D\ge40$
and a nonempty spare suffix.  At a current source instruction, the
unchanged decoder's established path semantics chooses one of the
listed leaf guard cases.  Its checked local simulation preserves the
dispatch invariant and implements the exact source successor.  The
span cannot decrease, and the spare suffix remains nonempty.  Induction
therefore simulates every finite source run.

If the source reaches its halt instruction, the corresponding checked
case reaches old state 254.  Conversely, suppose the source does not
halt.  Its successive instruction simulations form deterministic
positive-length machine segments.  Every segment is finite and has no
intermediate halt; otherwise it could not reach the certified target.
Their cumulative lengths tend to infinity, because each includes at
least one literal transition.  Thus every finite time after bootstrap
lies in a nonhalting segment.  The binary machine cannot halt.
\end{proof}

\begin{corollary}[Consistency endpoint]
Using the unchanged source-program and set-theoretic calibration, the
250-state machine halts exactly when ZFC has a formal contradiction.
Consequently the conventional proof gives
\[
 \mathrm{ZFC}\ \vdash
 \bigl(\text{the 250-state blank-tape machine does not halt}\bigr)
 \ \longleftrightarrow\ \Con(\mathrm{ZFC}).
\]
\end{corollary}

The new transport uses elementary finite-string and natural-number
induction.  The existing calibration is invoked as its ordinary
mathematical argument about the fixed numerical calculus, the SIR
set theory, ZF, and ZFC; its internal ZFC formulation is not inferred
merely from a Lean axiom report.  In particular, the usual model and
constructible-universe consistency arguments may be carried out inside
ZFC, with first-order completeness where needed.  A new proof-assistant
formalization of this exact 250-state endpoint, and an internalization
of the complete calibration in PA, remain separate work.

\section{Reproduction and attribution}

The accompanying scripts use the Python standard library:
\begin{verbatim}
python3 -B code/verify_all.py
\end{verbatim}
Run this command from the extracted archive root. It reconstructs both
frozen tables, checks the universally parameterized symbolic cases, runs
3,600 literal box regressions and the independent instruction and exact-time
audits, and checks the complete blank bootstrap.  The random
tests are explicitly secondary to the proofs above.  No search is
required to reproduce the final table or its certificates.

The original machine and compact calculus are due to Mira-Cult; Jeff
provided the original verification and symbolic rules, and Wade the
compiler framework.  The 253-state reduction is due to mira-emmy.
ProveIt supplied the permutation252 baseline and preserved the
attributed source-27 strict affine checker.  This contribution changes
the binary interpreter by shared prefix carry, rollback, spare
accounting, and pulse retiming.  Source identities are retained in
\texttt{provenance.json}; the unmodified source-27 checker is included
for comparison.

\begingroup\small\raggedright
\begin{thebibliography}{9}
\bibitem{candidate252}
ProveIt, \emph{A 252-state machine with a checked ZFC consistency endpoint},
\url{https://github.com/VladimirReshetnikov/ProveIt/tree/main/Computability/BusyBeaver/Research/ZFC/candidate-permutation252}.
\bibitem{m254}
Jeff, \emph{A verified 254-state Turing machine that halts iff ZFC is inconsistent},
including the attributed Mira-Cult table,
\url{https://github.com/jeff87654/busy-beaver/tree/zf-254/ZF/m254}.
\bibitem{tm-design}
Jeff, \emph{zfc\_254: Turing-machine layer},
\url{https://github.com/jeff87654/busy-beaver/blob/zf-254/ZF/m254/design/tm.md}.
\end{thebibliography}
\endgroup
\end{document}
